\begin{solution}
    Soit $z=a+ib\in\C$. Pour $t\in\Z$, on a $e^{itz}=e^{-bt}e^{ait}$.
    Deux cas se présentent :
    \begin{itemize}
        \item \textbf{Si $b \neq 0$ :} La suite des modules $|e^{itz}| = e^{-bt}$ tend vers $0$ ou $+\infty$ lorsque $t \to \pm\infty$.
        Pour tout compact $K \subset \C^*$, $K$ est borné et à distance strictement positive de $0$. L'intersection $K \cap G_z$ est donc nécessairement finie. $G_z$ est par conséquent une partie discrète de $\C^*$, ce qui implique qu'elle y est fermée.
        \item \textbf{Si $b = 0$ :} $z \in \R$. $G_z$ est un sous-groupe de $\mathbb U$.
        Si $z \in \pi\Q$, il existe $p, q \in \Z^*$ tels que $z = p\pi/q$. $G_z$ est alors le groupe fini des racines de l'unité d'ordre $2q$, qui est un fermé de $\C^*$.
        Si $z \notin \pi\Q$, on sait que $G_z$ est dense dans $\mathbb U$. Son adhérence dans $\C^*$ est donc $\mathbb U$, qui est strictement plus grand que $G_z$. $G_z$ n'est donc pas fermé.
    \end{itemize}
    En conclusion, $G_z$ est un fermé de $\C^*$ si et seulement si $z \notin \R \setminus \pi\Q$.
\end{solution}

\begin{solution}
    Soit $A\in X$. Considérons la suite $(A^n)_{n\in\N^*}$.
    Comme $A\in X$ et que $X$ est stable par produit, c'est une suite d'éléments de $X$. 
    $X$ étant compact, cette suite admet une valeur d'adhérence dans $X$. Il existe donc une extractrice $\varphi : \N\to\N$ et $B\in X$ tels que $A^{\varphi(n)}\longrightarrow B$.
    
    Remarquons que la stricte croissance de $\varphi$ impose $\varphi(n+1) - \varphi(n) \geq 1$. La suite $\left(A^{\varphi(n+1)-\varphi(n)}\right)_{n\in\N}$ est donc bien à valeurs dans $X$.
    Par continuité du produit et de l'inversion dans $\gl_n(\R)$, on a :
    \[ 
        A^{\varphi(n+1)-\varphi(n)} = A^{\varphi(n+1)} \left(A^{\varphi(n)}\right)^{-1} \longrightarrow B B^{-1} = I_n 
    \]
    Puisque $X$ est fermé, $I_n \in X$.
    
    Considérons maintenant la suite $M_n = A^{\varphi(n+1)-\varphi(n)-1}$. Puisque l'exposant est positif ou nul, $M_n \in X$ (avec $A^0 = I_n \in X$).
    De plus, on a la relation $M_n A = A^{\varphi(n+1)-\varphi(n)} \longrightarrow I_n$.
    $X$ étant compact, la suite $(M_n)$ admet une sous-suite convergente vers un certain $C \in X$.
    Par passage à la limite, $C A = I_n$. Dans $\gl_n(\R)$, un inverse à gauche est l'inverse de la matrice, donc $C = A^{-1}$. Ainsi, $A^{-1} \in X$.
    $X$ est donc stable par passage à l'inverse, c'est un sous-groupe.
\end{solution}
     
\begin{solution}
    \begin{enumerate}
        \item Montrons que $\R[X]\backslash\mathcal U$ est un ouvert de $(\R[X],\norm\cdot_\infty)$.
        Soit un polynôme $P \notin \mathcal U$. 
        Si $P = 0$, posons $\varepsilon = 1/2$. Tout polynôme $Q \in B(0, \varepsilon)$ a des coefficients strictement inférieurs à $1/2$, son coefficient dominant ne peut donc pas valoir $1$, d'où $Q \notin \mathcal U$.
        
        Si $P \neq 0$, notons $d = \deg P$ et $c = \cd P$. Comme $P \notin \mathcal U$, on a $c \neq 1$.
        Posons $\varepsilon = \frac{1}{2}\min(1, |c - 1|, |c|) > 0$.
        Soit $Q \in \mathcal U$. $Q$ possède un coefficient dominant égal à $1$ pour un certain degré $m$.
        \begin{itemize}
            \item Si $m > d$, le $m$-ème coefficient de $P$ est nul, donc $\norm{Q-P}_\infty \geq |1 - 0| = 1 > \varepsilon$.
            \item Si $m = d$, $\norm{Q-P}_\infty \geq |1 - c| > \varepsilon$.
            \item Si $m < d$, le $d$-ème coefficient de $Q$ est nul, donc $\norm{Q-P}_\infty \geq |0 - c| = |c| > \varepsilon$.
        \end{itemize}
        Dans tous les cas, la distance entre $P$ et n'importe quel polynôme de $\mathcal U$ est strictement supérieure à $\varepsilon$. Donc $B(P,\varepsilon) \subset \R[X]\backslash\mathcal U$, ce qui prouve que $\mathcal U$ est fermé.
        \item Supposons $Q$ scindé sur $\R$.
        Notons
        \[
            Q=\prod_{k=1}^p(X-\lambda_k),\ \lambda_k\in\R  
        \]
        Soit $z\in\C$. On a
        \[
            |Q(z)|^2=Q(z)\overline{Q(z)}=Q(z)Q(\overline z)=\prod_{k=1}^p(\lambda_k^2-2\lambda_k\Re(z)+|z|^2)
        \]
        On considère alors $y\in\R\mapsto y^2-2\Re(z)y+|z|^2$, une application polynomiale de degré $2$ atteignant son minimum en $\Re(z)$, et ce minimum vaut $|z|^2 - \Re(z)^2 = |\Im(z)|^2$.
        D'où, pour $k=1,\dots,p$, $\lambda_k^2-2\lambda_k\Re(z)+|z|^2\geq |\Im(z)|^2$, et donc finalement 
        \[
            |Q(z)|^2\geq\prod_{k=1}^p|\Im(z)|^2=|\Im(z)|^{2p} \implies |Q(z)|\geq|\Im(z)|^p
        \]
        Réciproquement, si pour tout $z\in \C$, $|Q(z)|\geq|\Im(z)|^p$, alors pour toute racine $\alpha\in\C$ de $Q$, on a $0=|Q(\alpha)|\geq|\Im(\alpha)|^p$, soit $\Im(\alpha)=0$ et $\alpha\in\R$.
        \item Soit $(Q_n)\in\mathcal S^{\N}$, et supposons que cette suite converge vers $P\in\R[X]$.
        Soit $p=\deg P$. On montre que l'on peut extraire de $(Q_n)$ une suite dont tous les termes sont des polynômes de degré égal à $p$.
        
        Pour cela, on procède par l'absurde en supposant qu'à partir d'un certain rang, tous les $Q_n$ sont de degré $\neq p$.
        Si une infinité de $Q_n$ ont un degré $> p$, il existe une extractrice $\varphi$ telle que $\deg Q_{\varphi(n)} \ge p+1$. 
        Puisque $Q_{\varphi(n)}$ est unitaire, son coefficient de degré $\deg Q_{\varphi(n)}$ vaut $1$. Or, ce coefficient pour $P$ vaut $0$.
        Donc $\norm{Q_{\varphi(n)} - P}_\infty \ge 1$, ce qui contredit la convergence vers $P$ (une contradiction similaire apparaît sur le coefficient de degré $p$ si une infinité de termes ont un degré $< p$).
        
        On dispose donc d'une extractrice $\varphi:\N\to\N$ telle que pour tout $n\in\N$, $\deg Q_{\varphi(n)}=p$.
        On a alors 
        \[
            \forall n\in\N,\forall z\in\C,\ |Q_{\varphi(n)}(z)|\geq|\Im(z)|^p    \tag*{(*)}
        \]
        Comme $Q_{\varphi(n)} \to P$ au sens de la norme $\norm{\cdot}_\infty$, on a $Q_{\varphi(n)}(z) \to P(z)$ pour tout $z \in \C$. 
        En passant à la limite dans (*), on obtient
        \[
            \forall z\in\C,\ |P(z)|\geq|\Im(z)|^p
        \]
        ce qui assure que $P$ est scindé sur $\R$. La première question assure de plus que $P$ est unitaire. $\mathcal S$ est donc bien un fermé de $\R[X]$.
    \end{enumerate}
\end{solution}